\documentclass[10pt,a4paper]{amsart}
\usepackage{amsmath,amssymb,amsthm,graphicx}
% For dealing with references we use the comment environment
\usepackage{verbatim}
\newenvironment{reference}{\comment}{\endcomment}
%\newenvironment{reference}{}{}
\newenvironment{slogan}{\comment}{\endcomment}
\newenvironment{history}{\comment}{\endcomment}

% For commutative diagrams we use Xy-pic
\usepackage[all]{xy}
% Theorem environments.
%
\theoremstyle{plain}
\newtheorem{theorem}[subsection]{Theorem}
\newtheorem{proposition}[subsection]{Proposition}
\newtheorem{lemma}[subsection]{Lemma}

\theoremstyle{definition}
\newtheorem{definition}[subsection]{Definition}
\newtheorem{example}[subsection]{Example}
\newtheorem{exercise}[subsection]{Exercise}
\newtheorem{situation}[subsection]{Situation}

\theoremstyle{remark}
\newtheorem{remark}[subsection]{Remark}
\newtheorem{remarks}[subsection]{Remarks}

\numberwithin{equation}{subsection}

% Macros
%
\def\lim{\mathop{\mathrm{lim}}\nolimits}
\def\colim{\mathop{\mathrm{colim}}\nolimits}
\def\Spec{\mathop{\mathrm{Spec}}}
\def\Hom{\mathop{\mathrm{Hom}}\nolimits}
\def\Ext{\mathop{\mathrm{Ext}}\nolimits}
\def\SheafHom{\mathop{\mathcal{H}\!\mathit{om}}\nolimits}
\def\SheafExt{\mathop{\mathcal{E}\!\mathit{xt}}\nolimits}
\def\Sch{\mathit{Sch}}
\def\Mor{\mathop{\mathrm{Mor}}\nolimits}
\def\Ob{\mathop{\mathrm{Ob}}\nolimits}
\def\Sh{\mathop{\mathit{Sh}}\nolimits}
\def\NL{\mathop{N\!L}\nolimits}
\def\CH{\mathop{\mathrm{CH}}\nolimits}
\def\proetale{{pro\text{-}\acute{e}tale}}
\def\etale{{\acute{e}tale}}
\def\QCoh{\mathit{QCoh}}
\def\Ker{\mathop{\mathrm{Ker}}}
\def\Im{\mathop{\mathrm{Im}}}
\def\Coker{\mathop{\mathrm{Coker}}}
\def\Coim{\mathop{\mathrm{Coim}}}

% Boxtimes
%
\DeclareMathSymbol{\boxtimes}{\mathbin}{AMSa}{"02}

%
% Macros for moduli stacks/spaces
%
\def\QCohstack{\mathcal{QC}\!\mathit{oh}}
\def\Cohstack{\mathcal{C}\!\mathit{oh}}
\usepackage{cleveref}
\crefname{originallemma}{original lemma}{original lemmas}
\Crefname{originallemma}{Original lemma}{Original lemmas}
\makeatletter
\newlabel{lemma-subcategory-right-resolution}{{05T6}{source}{}{sourcepointer-1}{}}
\newlabel{lemma-subcategory-right-resolution@cref}{{[originallemma][0][]05T6}{[1][0][]source}{}{}{}}
\makeatother
\title{Locally finite elementary replacements produce a dimension-zero resolution of an unbounded complex}
\author{The Stacks Project Authors}
\date{stacks-project:derived.tex:local-snapshot-20261011; GNU FDL 1.2 or later}
\begin{document}
\maketitle
\noindent\textit{Editorial note.} Before the proof starts, reindex by a finite shift so every term with d(K\textasciicircum{}n)=infinity lies in degree>=0. Such a shift exists by n+d(K\textasciicircum{}n)→-infinity as n→-infinity. The initial bounded-below replacement then removes these terms before \ensuremath{\xi} is asserted finite. Original proof wording is retained.\par\medskip
\noindent\textit{Editorial scope.} One original target with its local construction proved in full. Classical category primitives are explicitly allowed.\par\medskip
\noindent\textit{Allowed standard prerequisites (editorial).}\begin{itemize}
\item Bounded-below resolutions by a monomorphism-cogenerating class. A bounded-below complex embeds quasi-isomorphically in a bounded-below complex whose terms lie in a given class containing zero and into which every object embeds.
\item Elementary abelian-category pushouts and local stabilization. A pushout replacement has an acyclic quotient; if each fixed finite-degree window eventually stops changing, cohomology is preserved in the stabilized complex.
\end{itemize}\medskip
\ifdefined\hypertarget\hypertarget{sourcepointer-1}{}\fi
\noindent\textit{Source reference.} Original lemma 05T6: focused-seeds/ACL/sources/stacks-snapshot-20261011/derived.tex, source lines 5386--5401. Standard bounded-below replacement used only to initialize; the locally finite unbounded schedule is not given. Stable Stacks tag displayed.\par\medskip
\begin{lemma}
\label{lemma-replace-resolution}
Let $\mathcal{A}$ be an abelian category. Let
$d : \Ob(\mathcal{A}) \to \{0, 1, 2, \ldots, \infty\}$ be a function.
Assume that
\begin{enumerate}
\item every object of $\mathcal{A}$ is a subobject of an
object $A$ with $d(A) = 0$,
\item $d(A \oplus B) \leq \max \{d(A), d(B)\}$ for $A, B \in \mathcal{A}$, and
\item if $0 \to A \to B \to C \to 0$ is short exact, then
$d(C) \leq \max\{d(A) - 1, d(B)\}$.
\end{enumerate}
Let $K^\bullet$ be a complex such that $n + d(K^n)$ tends to $-\infty$
as $n \to -\infty$. Then there exists a quasi-isomorphism
$K^\bullet \to L^\bullet$ with $d(L^n) = 0$ for all $n \in \mathbf{Z}$.
\end{lemma}

\begin{proof}
By Lemma \ref{lemma-subcategory-right-resolution} we can find a
quasi-isomorphism $\sigma_{\geq 0}K^\bullet \to M^\bullet$ with
$M^n = 0$ for $n < 0$ and $d(M^n) = 0$ for $n \geq 0$. Then $K^\bullet$
is quasi-isomorphic to the complex
$$
\ldots \to K^{-2} \to K^{-1} \to M^0 \to M^1 \to \ldots
$$
Hence we may assume that $d(K^n) = 0$ for $n \gg 0$. Note that
the condition $n + d(K^n) \to -\infty$ as $n \to -\infty$ is not
violated by this replacement.

\medskip\noindent
We are going to improve $K^\bullet$ by an (infinite) sequence of
elementary replacements. An {\it elementary replacement} is the following.
Choose an index $n$ such that $d(K^n) > 0$. Choose an injection
$K^n \to M$ where $d(M) = 0$. Set
$M' = \Coker(K^n \to M \oplus K^{n + 1})$. Consider the map of complexes
$$
\xymatrix{
K^\bullet : \ar[d] &
K^{n - 1} \ar[d] \ar[r] &
K^n \ar[d] \ar[r] &
K^{n + 1} \ar[d] \ar[r] &
K^{n + 2} \ar[d] \\
(K')^\bullet : &
K^{n - 1} \ar[r] &
M \ar[r] &
M' \ar[r] &
K^{n + 2}
}
$$
It is clear that $K^\bullet \to (K')^\bullet$ is a quasi-isomorphism.
Moreover, it is clear that $d((K')^n) = 0$ and
$$
d((K')^{n + 1}) \leq \max\{d(K^n) - 1, d(M \oplus K^{n + 1})\} \leq
\max\{d(K^n) - 1, d(K^{n + 1})\}
$$
and the other values are unchanged.

\medskip\noindent
To finish the proof we carefully choose the order in which to do
the elementary replacements so that for every integer $m$ the complex
$\sigma_{\geq m}K^\bullet$ is changed only a finite number of times.
To do this set
$$
\xi(K^\bullet) = \max \{n + d(K^n) \mid d(K^n) > 0\}
$$
and
$$
I = \{n \in \mathbf{Z} \mid \xi(K^\bullet) = n + d(K^n)
\text{ and }
 d(K^n) > 0\}
$$
Our assumption that $n + d(K^n)$ tends to $-\infty$ as $n \to -\infty$
and the fact that $d(K^n) = 0$ for $n >> 0$
implies $\xi(K^\bullet) < +\infty$ and that $I$ is a finite set.
It is clear that $\xi((K')^\bullet) \leq \xi(K^\bullet)$ for an
elementary transformation as above. An elementary transformation
changes the complex in degrees $\leq \xi(K^\bullet) + 1$. Hence if we can
find finite sequence of elementary transformations which
decrease $\xi(K^\bullet)$, then we win. However, note that if we
do an elementary transformation starting with the smallest element
$n \in I$, then we either decrease the size of $I$, or we increase
$\min I$. Since every element of $I$ is $\leq \xi(K^\bullet)$ we see
that we win after a finite number of steps.
\end{proof}

\end{document}
